\section{Complementary divisors and affine parity fibers}
\label{sec:complement-fibers}

The dynamic theorem stops before the Poisson-short or conductor margin
is exhausted.  We now give an exact normal form for what remains.  No
estimate in this section assumes that a target is squarefree.

\subsection{Hyperbola reflection}

For \(n\ge1\) and \(Y\ge1\), put
\begin{equation}
 \cT_Y(n)=\sum_{\substack{u\mid n\\u>Y}}a(u),
 \qquad a(u)=-\mu(u)\log u.
 \label{eq:TY}
\end{equation}

\begin{proposition}[Exact complement and Liouville cocycle]
\label{prop:complement-cocycle}
For every \(n,Y\ge1\),
\begin{align}
 \cT_Y(n)
 &=\sum_{\substack{r\mid n\\r<n/Y}}a(n/r)
 \label{eq:exact-complement}\\
 &=-\lambda(n)
 \sum_{\substack{r\mid n\\r<n/Y}}
 \lambda(r)\mu^2(n/r)\log(n/r).
 \label{eq:exact-Liouville-cocycle}
\end{align}
Here \(\lambda\) is the Liouville function.  In particular, the sign on
the reflected quotient is \(\mu(n/r)\), not \(\mu(r)\).
\end{proposition}

\begin{proof}
The involution \(u\mapsto r=n/u\) bijects divisors satisfying \(u>Y\)
with those satisfying \(r<n/Y\), proving
\eqref{eq:exact-complement}.  Since
\(\mu(k)=\lambda(k)\mu^2(k)\) and \(\lambda\) is completely
multiplicative,
\[
 \mu(n/r)=\lambda(n/r)\mu^2(n/r)
 =\lambda(n)\lambda(r)\mu^2(n/r).
\]
Insert this identity into \(a(n/r)\).
\end{proof}

If \(n>Y^2\), the tail has the exact hyperbola split
\begin{equation}
 \cT_Y(n)
 =\underbrace{\sum_{\substack{r\mid n\\r\le Y}}a(n/r)}_{\cC_Y(n)}
 +\underbrace{\sum_{\substack{r\mid n\\Y<r<n/Y}}a(n/r)}_{\cB_Y(n)},
 \label{eq:balanced-split}
\end{equation}
with the evident adjustment when a boundary divisor occurs.  The first
piece has a short complementary divisor.  In the second, both factors
of \(n=r(n/r)\) exceed \(Y\); it is a genuine Type--II hyperbola band.
Pairing divisors about \(\sqrt n\) gives
\begin{align}
 \cB_Y(n)
 ={}&\sum_{\substack{d\mid n\\Y<d<\sqrt n}}
 \{a(d)+a(n/d)\}
 +\one_{n=z^2,\ z>Y}a(z).
 \label{eq:folded-balanced-band}
\end{align}
On a squarefree target, the paired coefficient is
\begin{equation}
 a(d)+a(n/d)
 =-\lambda(d)\{\log d+\lambda(n)\log(n/d)\}.
 \label{eq:squarefree-parity-pair}
\end{equation}
Thus even and odd factor parity are different eigensectors of divisor
reflection.  This is an identity, not a bound for either sector.

For the TPC target \(n\asymp X\) and
\(Y=R=X^{1/2-\delta}\), the balanced factor exponents lie between
\(1/2-\delta\) and \(1/2+\delta\).  The dynamic theorem removes a
strict lower part of this geometry before the quotient-M\"obius
normal form is needed.

\subsection{Two complement fibers}

Fix distinct positive primitive rows \(m_1,m_2\), meaning
\((m_1m_2,h)=1\), and write
\begin{equation}
 n_i(j)=m_i j+h\qquad(i=1,2).
 \label{eq:two-targets}
\end{equation}
Let \(r,s\ge1\).  Since \((m_1m_2,h)=1\), any solution of
\(r\mid n_1(j)\), \(s\mid n_2(j)\) necessarily satisfies
\begin{equation}
 (r,m_1)=(s,m_2)=1.
 \label{eq:complement-primitivity}
\end{equation}
We restrict to such pairs and put
\begin{equation}
 d=(r,s),\qquad q=[r,s]=\frac{rs}{d}.
 \label{eq:rsdq}
\end{equation}

\begin{theorem}[Exact affine complement fiber]
\label{thm:affine-complement-fiber}
The two congruences
\begin{equation}
 r\mid m_1j+h,
 \qquad s\mid m_2j+h
 \label{eq:two-complement-congruences}
\end{equation}
are jointly soluble if and only if
\begin{equation}
 d\mid h(m_1-m_2).
 \label{eq:complement-compatibility}
\end{equation}
When they are soluble, they select a unique class \(j_0\pmod q\).
Writing \(j=j_0+qt\), define
\begin{align}
 U_{r,s}(t)
 &=\frac{m_1(j_0+qt)+h}{r}=A_1t+B_1,
 \label{eq:U-form}\\
 V_{r,s}(t)
 &=\frac{m_2(j_0+qt)+h}{s}=A_2t+B_2.
 \label{eq:V-form}
\end{align}
These are integer affine forms with
\begin{equation}
 \boxed{
 A_1B_2-A_2B_1
 =\frac{h(m_1-m_2)}d\ne0.}
 \label{eq:affine-determinant}
\end{equation}
Consequently, for any finitely supported orbit weight \(W\),
\begin{align}
 &\sum_jW(j)\cT_S(n_1(j))\cT_S(n_2(j))
 \notag\\
 &\quad=
 \sum_{\substack{r,s\ge1\\(r,m_1)=(s,m_2)=1\\
                  (r,s)\mid h(m_1-m_2)}}
 \ \sum_{t\in\Z}
 W(j_0+qt)
 a(U_{r,s}(t))a(V_{r,s}(t))
 \one_{U_{r,s}(t)>S}
 \one_{V_{r,s}(t)>S}.
 \label{eq:exact-two-fiber-reassembly}
\end{align}
All terms with nonpositive targets or outside the finite horizon are
understood to be absent.
\end{theorem}

\begin{proof}
Each congruence in \eqref{eq:two-complement-congruences} selects one
class modulo its divisor.  Their reductions modulo \(d\) agree exactly
when
\(-h\overline{m_1}\equiv-h\overline{m_2}\pmod d\).  Since
\((m_1m_2,d)=1\), multiplication by \(m_1m_2\) shows that this is
equivalent to
\eqref{eq:complement-compatibility}.  The Chinese remainder theorem
then gives one class modulo \(q\).

The quantities in \eqref{eq:U-form}--\eqref{eq:V-form} are integers,
with
\[
 A_1=\frac{m_1q}{r},\quad
 B_1=\frac{m_1j_0+h}{r},\quad
 A_2=\frac{m_2q}{s},\quad
 B_2=\frac{m_2j_0+h}{s}.
\]
Direct calculation gives
\[
 A_1B_2-A_2B_1
 =\frac{qh(m_1-m_2)}{rs}
 =\frac{h(m_1-m_2)}d.
\]
It is nonzero because the generic mask has \(m_1\ne m_2\).
Finally, expand both tails using \eqref{eq:exact-complement}, group by
\((r,s)\), and parameterize the common progression.  The inequalities
on \(U,V\) are exactly the original requirements that the quotient
divisors exceed \(S\).
\end{proof}

Because \(a(U)a(V)=\mu(U)\mu(V)\log U\log V\),
\eqref{eq:exact-two-fiber-reassembly} is an exact family of two-point
affine M\"obius correlations.  The determinant content \(d=(r,s)\)
divides \(h(m_1-m_2)\).  On the primitive orbit \((j,H)=1\), it is
coprime to \(H\) and therefore divides the row difference.  In a mixed
large--short leg there is also
the original divisor gcd \((U,v)\mid m_1-m_2\); the two contents can
overlap when the target has square factors.  They must not be declared
coprime.

\begin{remark}[The \(r=s=1\) parity core]
\label{rem:unit-complement}
The first complement fiber is always compatible and has
\begin{equation}
 U_{1,1}(j)=m_1j+h,\qquad
 V_{1,1}(j)=m_2j+h.
 \label{eq:unit-complement-forms}
\end{equation}
Its coefficient is
\begin{equation}
 \mu(m_1j+h)\mu(m_2j+h)
 \log(m_1j+h)\log(m_2j+h).
 \label{eq:unit-parity-coefficient}
\end{equation}
On prime targets, each individual factor
\(-\mu(m_i j+h)\log(m_i j+h)\) is the positive prime weight.  Thus a
generic assertion that this fiber is negligible by ``M\"obius
randomness'' would be a parity-level overclaim; the precise averaging
and uniformity are indispensable.
\end{remark}
